\documentclass[12pt]{article}
\usepackage[utf8]{inputenc}
\usepackage[spanish]{babel}
\usepackage{amsmath}
\usepackage{amssymb}
\usepackage{geometry}
\geometry{a4paper, margin=1in}
\begin{document}
\section{Módulo: src/core (Motor Numérico y Físico)}

El directorio \texttt{src/core} constituye el cerebro matemático de la simulación. Su responsabilidad es pura y exclusivamente el cálculo numérico, la dinámica geométrica y la actualización de los estados físicos del universo, permaneciendo completamente ciego respecto a cómo estos datos se dibujan en pantalla o cómo suenan.

\subsection{Archivos Principales y Funcionalidad}
\begin{itemize}
    \item \textbf{\texttt{nucleo\_espectral.py} y \textbf{\texttt{nucleo\_neural.py}}:} Estos archivos contienen las clases matemáticas que resuelven las ecuaciones diferenciales (como Gray-Scott o Navier-Stokes) y las redes neuronales (CPPN). Manejan tensores de NumPy/PyTorch para calcular gradientes espaciales y operadores laplacianos en cada iteración temporal.
    \item \textbf{\texttt{visual\_entities.py}:} Define la orientación espacial, vectores de posición, velocidad y aceleración de las partículas discretas en simulaciones de N-cuerpos. 
    \item \textbf{\texttt{efectos\_visuales.py}:} Gestiona las transformaciones proyectivas de la matriz espacial, tales como las rotaciones en ejes 3D proyectadas a 2D y el factor de escala (zoom) dictaminado por la intensidad física calculada.
\end{itemize}

\subsection{Cálculos de Dinámica y Cinemática}
En este módulo, el tiempo se trata como una variable discreta. En cada actualización (\textit{update}), las coordenadas y matrices de estado evolucionan. Se resuelven integraciones explícitas de Euler o Runge-Kutta para asegurar que la masa, el momento y la energía matemática se comporten bajo leyes deterministas. No existen colores ni píxeles aquí, sólo arreglos de números flotantes de alta precisión.

\end{document}
